\documentclass[a4paper,11pt]{ltjsarticle}
\usepackage[margin=22mm]{geometry}
\usepackage{amsmath}
\usepackage{booktabs}
\usepackage{listings}
\usepackage{xcolor}
\usepackage{hyperref}
\hypersetup{colorlinks=true,linkcolor=blue,urlcolor=blue,citecolor=blue}

\lstset{
  basicstyle=\ttfamily\footnotesize,
  backgroundcolor=\color{black!4},
  frame=single,
  framerule=0pt,
  breaklines=true,
  showstringspaces=false,
  numbers=left,
  numberstyle=\tiny\color{gray},
  numbersep=6pt,
  xleftmargin=12pt,
  keywordstyle=\color{blue!70!black}\bfseries,
  commentstyle=\color{green!40!black}\itshape,
  stringstyle=\color{red!60!black},
  language=C
}

\title{\textbf{Embedded Xinu 上での Raspberry Pi 3 オンボード WiFi\\(BCM43455)ドライバのゼロからの実装}\\[4pt]
\large ---- スキャンから WPA2 接続・DHCP・ping 疎通まで ----}
\author{実機検証レポート（arm-rpi3 platform）}
\date{2026年6月2日}

\begin{document}
\maketitle

\begin{abstract}
本レポートは、教育用 OS である Embedded Xinu を Raspberry Pi 3 B+（BCM2837）へ移植した
\texttt{arm-rpi3} プラットフォーム上に、オンボード WiFi チップ Broadcom/Cypress BCM43455 の
ドライバを\textbf{ゼロから実装}し、アクセスポイントのスキャン・一覧表示から WPA2-PSK 接続、
DHCP による IP 取得、そして ping 疎通までを実機で達成した記録である。Linux の \texttt{brcmfmac}
ドライバおよび plan9/circle の \texttt{ether4330} ドライバをリファレンスとしつつ、Xinu の限られた
環境（MMU 有効・D-cache 無効・標準 C ライブラリ僅少）に合わせて SDIO バスの bring-up、ファーム
ウェアのロード、SDPCM/BDC フレーミング、ファームウェア内蔵サプリカント(FWSUP)による 4-way
ハンドシェイク、最小限の IP スタック(ARP/ICMP)を実装した。本実装で突破した\textbf{6 つの技術的な壁}
と、その診断・解決の過程を詳述する。最終的に \texttt{ping 192.168.3.52} は定常状態で
\textbf{パケットロス 0\%・平均 RTT 15\,ms} を達成した。
\end{abstract}

\section{背景と目標}
Embedded Xinu は小規模で読みやすい教育用 OS であり、本プロジェクトではこれを Raspberry Pi 3 B+
の実機（BCM2837、Cortex-A53、32-bit）へ移植している。Pi 3 B+ のオンボード WiFi は BCM43455 で、
SoC とは SDIO バス（Arasan SDHCI/EMMC コントローラ、物理アドレス \texttt{0x3F300000}）で接続されている。

目標は単純である：\textbf{「WiFi アクセスポイントの一覧を表示し、その中から選んで接続する」}。
ただし Xinu には WiFi スタックが一切存在しないため、SDIO の最下層からファームウェアのロード、
無線制御プロトコル、暗号化(WPA2)、上位の IP 疎通までを自前で実装する必要があった。

\section{システム構成}
実装は単一の自己完結したファイル \texttt{apps/wifi.c}（約 2{,}000 行）に集約した。
HTTP 経由の操作・診断インタフェースは \texttt{apps/webactor.c} が提供する。レイヤ構成を表
\ref{tab:layers} に示す。

\begin{table}[h]
\centering
\caption{WiFi ドライバのレイヤ構成}
\label{tab:layers}
\small
\begin{tabular}{@{}lll@{}}
\toprule
レイヤ & 内容 & 主な関数 \\
\midrule
SDIO host & Arasan SDHCI 初期化, CMD52/53 & \texttt{emmc\_*}, \texttt{sdio\_cmd52/53} \\
バックプレーン & EROM コアスキャン, ARM CR4 制御 & \texttt{wifi\_sbmem}, \texttt{wifi\_corescan} \\
firmware ロード & fw/nvram/clm を RAM へ転送 & \texttt{wifi\_fwload}, \texttt{wifi\_clmload} \\
SDPCM/BDC & 制御(ioctl)・データのフレーミング & \texttt{wifi\_wlcmd}, \texttt{wifi\_data\_tx} \\
無線制御 & scan, join, 暗号設定 & \texttt{wifi\_scan}, \texttt{wifi\_do\_join} \\
IP 応答 & DHCP クライアント, ARP/ICMP 応答 & \texttt{wifi\_dhcp}, \texttt{wifi\_net\_service} \\
\bottomrule
\end{tabular}
\end{table}

ファームウェア（\texttt{cyfmac43455-sdio-minimal.bin}, 約 548\,KB）、NVRAM、CLM(規制)ブロブは
\texttt{.incbin} でカーネルイメージに直接埋め込み、SDIO 経由でチップ RAM(\texttt{0x198000})へ
ストリーム転送する。MMU は有効化しているが D-cache は USB-ethernet/フレームバッファの DMA を
壊すため意図的に無効、命令キャッシュ(I-cache)のみ有効化している。

\section{実装の道のり：6 つの壁}
スキャン（AP 一覧の取得）までは比較的順調に到達したが、\textbf{WPA2 接続から先}には
段階的に 6 つの壁が立ちはだかった。各々を実機ログによる診断で特定し、解決した（表 \ref{tab:walls}）。

\begin{table}[h]
\centering
\caption{突破した 6 つの技術的な壁}
\label{tab:walls}
\small
\begin{tabular}{@{}clp{8.4cm}@{}}
\toprule
\# & 症状 & 原因と解決 \\
\midrule
1 & \texttt{sup\_wpa} が \texttt{-23} & 使用中の \emph{standard} ファームウェアに在荷サプリカント(FWSUP)が無い。
   \texttt{idsup} を持つ \emph{minimal} ファームウェア(7.45.241)へ差し替え。\\
2 & \texttt{WSEC\_PMK} が \texttt{-2} & PMK 構造体 \texttt{brcmf\_wsec\_pmk\_le} は \texttt{key[128]} を含む
   \textbf{132 バイト固定}。36 バイトで送っていたため拒否されていた。\\
3 & 接続しても assoc せず & \texttt{wsec}/\texttt{wpa\_auth}/\texttt{auth} をコマンド(WLC\_SET\_*)で設定していたが、
   join/FWSUP が読むのは \textbf{bsscfg iovar}。全て iovar 設定に変更。\\
4 & join が無反応(scan のみ) & ext\_join 構造体 \texttt{brcmf\_ext\_join\_params\_le} は \texttt{\_\_packed} ではなく
   \textbf{自然アラインメントで 70 バイト}。packed の 65 バイトは \texttt{-14}(BUFTOOSHORT)、
   114 バイトはフィールドがズレて誤読されていた。\\
5 & ping 不通(0\% 応答) & SDPCM の\textbf{送信クレジット window/フロー制御を無視}。fw は各フレームに
   \texttt{max\_seq}(byte 9)と \texttt{fcmask}(byte 8)を載せ、window 超過の送信は黙って破棄する。
   RX から追従し、開くまで待って送信。\\
6 & 接続が不安定 & association が AUTH で停滞し再スキャンに戻ることがある。join を最大 4 回リトライ。\\
\bottomrule
\end{tabular}
\end{table}

\subsection{壁 1・2：FWSUP ファームウェアと PMK 構造体}
WPA2-PSK の 4-way ハンドシェイクは、ホスト側サプリカントを書く代わりに、ファームウェア内蔵
サプリカント(FWSUP)に任せる方針とした。\texttt{brcmfmac} は \texttt{sup\_wpa} iovar の GET が
\texttt{-23}(UNSUPPORTED)を返さないことで FWSUP を判定する（\texttt{feature.c:349}）。RPi 配布の
\emph{standard} ファームウェアはこれを返した。capability 文字列を比較すると、\emph{minimal}
ファームウェアのみ \texttt{idsup-idauth}(在荷サプリカント)を持ち、\emph{standard} は代わりに
external-SAE(\texttt{extsae-dpp-sr-okc})を採用していた。\emph{minimal} へ差し替えると
\texttt{sup\_wpa=1} が \texttt{rc=0} で通り、PMK 構造体を 132 バイトに直すと \texttt{WSEC\_PMK} も
\texttt{rc=0} となった。

\subsection{壁 4：ext\_join 構造体のアラインメント（最大の難所）}
最も解決に時間を要したのがこの壁である。\texttt{join} iovar に渡す構造体について、当初は
\texttt{brcmf\_scan\_params\_le}(入れ子の SSID を含む大きな構造体)を仮定して 114 バイトで送って
いたが、ファームウェアは長さは受理するものの scan/assoc フィールドを誤読し、probe だけして
association に進まなかった。新しい \texttt{brcmfmac} の \texttt{brcmf\_join\_scan\_params\_le} を
仮定して 65 バイトに縮めると今度は \texttt{-14}(BUFTOOSHORT)。

決め手は構造体の\textbf{パッキングの再確認}であった。\texttt{fwil\_types.h} 全体で \texttt{\_\_packed}
は 1 箇所のみで、join 関連の構造体には付いていない。すなわち\textbf{自然アラインメント}であり、
実サイズは
\[
  \underbrace{\mathtt{offsetof(assoc\_le)}}_{56}
  + \underbrace{\mathtt{offsetof(chanspec\_list)}}_{12}
  + \underbrace{\mathtt{sizeof(u16)}}_{2}
  = 70\ \text{バイト}
\]
である。これは plan9 の \texttt{ether4330} が使う \texttt{wl\_extjoin\_params\_t} とバイト単位で一致した
（\texttt{scan\_type@36}, \texttt{bssid@56}, \texttt{chanspec@68}）。70 バイトの正しいレイアウトに直すと、
スキャンで得た実 chanspec を指定した directed join で初めて association が成立した。

\begin{lstlisting}[caption={壁 4 を解いた 70 バイト join 構造体（抜粋）}]
/* wl_extjoin_params_t (自然アラインメントで 70B):
 *   [0..3] SSID_len  [4..35] SSID[32]
 *   [36] scan_type(+3pad)  [40] nprobes  [44] active  [48] passive  [52] home
 *   [56..61] assoc bssid  [62..63] pad  [64..67] chanspec_num  [68..69] chanspec */
jp[0] = sl;                                  /* SSID_len */
for (i = 0; i < sl; i++) jp[4 + i] = ssid[i];
jp[36] = 0xFF;                               /* scan_type */
jp[40] = 2;                                  /* nprobes = 2 */
jp[44] = 120;                                /* active_time = 120 ms */
jp[48] = 0x86; jp[49] = 0x01;                /* passive_time = 390 ms */
jp[52]=jp[53]=jp[54]=jp[55]=0xFF;            /* home_time = -1 */
for (i = 0; i < 6; i++) jp[56 + i] = bssid[i];   /* assoc bssid */
jp[64] = 1;                                  /* chanspec_num = 1 */
jp[68] = chanspec & 0xFF; jp[69] = chanspec >> 8;
jsz = 70;
\end{lstlisting}

\subsection{壁 5：SDPCM 送信クレジットとフロー制御（ping 不通の真因）}
association・DHCP まで成功した後、\texttt{ping} が 100\% ロスした。受信フレームのログを追加して
診断したところ、Pi は\textbf{ARP 要求も ICMP echo 要求も正しく受信し、応答を生成・送出していた}
にもかかわらず、ホストに届いていなかった。原因は SDPCM の送信フロー制御を無視していたことに
あった。ファームウェアは送出する全フレームの SDPCM ヘッダに送信クレジット窓
\texttt{max\_seq}(byte 9)とチャネル別フロー制御マスク \texttt{fcmask}(byte 8)を載せており、
\texttt{txseq == txwindow} の状態やデータチャネルの FC が立った状態で送信したフレームは
\textbf{黙って破棄}される。早期に送った ARP 応答はたまたま窓内で届いたが、後続の ICMP 応答は
窓を超過して破棄されていた。受信フレームから窓とマスクを常時追従し、窓が開くまで待ってから
送信するように修正して解決した。

\section{接続シーケンスと実機検証}
最終的な接続シーケンスを以下に示す。
\begin{enumerate}
  \item \texttt{/api/wifi/scan} ---- escan により AP を列挙し JSON で返す
  \item \texttt{/api/wifi/join?ssid=...\&pass=...} ---- DOWN/INFRA/UP $\to$ scan で対象の BSSID/chanspec を特定
        $\to$ \texttt{wpa\_auth/auth/wsec} を iovar 設定 $\to$ \texttt{sup\_wpa=1} $\to$ \texttt{WSEC\_PMK}(PMK)
        $\to$ 70B \texttt{join} iovar $\to$ ファームウェアが 4-way を実行 $\to$ \texttt{E\_LINK} up
  \item \texttt{/api/wifi/dhcp} ---- DISCOVER/OFFER/REQUEST/ACK で IP を取得し、ARP/ICMP 応答スレッドを起動
\end{enumerate}

実機での join イベント列（\texttt{X-Wifi-EvSeq} ヘッダ）は
\texttt{124(scan), 46(PSK\_SUP, status 6=KEYED), 1(JOIN), 0(SET\_SSID)} となり、ファームウェア
サプリカントによる 4-way ハンドシェイク完了が確認できた。\texttt{GET\_BSSID} は AP の BSSID
\texttt{0c:73:29:8e:f8:2c} を返し association を裏付けた。DHCP は
\texttt{IP=192.168.3.52, mask=255.255.255.0, gw=192.168.3.1, dns=192.168.3.1} を取得した。

\begin{table}[h]
\centering
\caption{実機 ping 結果（WiFi 経由、定常状態 10 パケット）}
\label{tab:ping}
\small
\begin{tabular}{@{}lr@{}}
\toprule
指標 & 値 \\
\midrule
送信 / 受信 & 10 / 10 \\
パケットロス & 0.0\,\% \\
RTT 最小 / 平均 / 最大 & 11.5 / 15.3 / 28.0\,ms \\
\bottomrule
\end{tabular}
\end{table}

\begin{lstlisting}[numbers=none,caption={実機 ping 出力（抜粋）}]
64 bytes from 192.168.3.52: icmp_seq=1 ttl=64 time=11.533 ms
64 bytes from 192.168.3.52: icmp_seq=2 ttl=64 time=14.122 ms
 ... (10/10) ...
10 packets transmitted, 10 packets received, 0.0% packet loss
round-trip min/avg/max/stddev = 11.533/15.338/27.991/4.674 ms
\end{lstlisting}

初回試行では最初の数パケットが \texttt{Host is down}（ARP 解決中の settling）で落ちるが、
解決後は 0\% ロスで安定する。

\section{診断手法上の工夫}
長時間(約 150 秒)の join 処理の後では Xinu の TCP が大きな応答ボディを配送できず(write がブロック)、
従来の「全ログを HTTP ボディで返す」方式が機能しなかった。そこで、(1)主要な診断値を
\texttt{X-Wifi-*} レスポンス\textbf{ヘッダ}に載せる（小さなヘッダは確実に届く）、(2)即時応答する独立
エンドポイント \texttt{/api/wifi/trace} で完全ログを別取得する、という二段構えで診断を進めた。
イベント列を \texttt{X-Wifi-EvSeq} に CSV で載せたことが、壁 4・5 の特定に決定的に寄与した。

\section{結論と残課題}
Embedded Xinu 上に BCM43455 ドライバをゼロから実装し、\textbf{スキャン $\to$ AP 選択 $\to$
WPA2-PSK 接続 $\to$ DHCP $\to$ ping 疎通}までを実機で達成した。鍵となったのは、(i)在荷サプリカント
付きファームウェアの選択、(ii)ワイヤ構造体のサイズ・アラインメントの厳密な一致、(iii)SDPCM 送信
クレジット制御の遵守、の 3 点であり、いずれも実機ログによる地道な診断で特定した。

残課題として、TCP を含む上位プロトコル、再接続・ローミング、複数 AP の永続設定、
スキャン一覧からブラウザで選択・接続する UI などが挙げられる。

\vspace{4mm}
\noindent\footnotesize
リファレンス：Linux \texttt{brcmfmac}（\texttt{drivers/net/wireless/broadcom/brcm80211}）、
plan9/circle \texttt{ether4330}。ビルド：\texttt{arm-rpi3-port} ブランチ、コミット \texttt{2c4f783}。
デプロイは SD カードの差し替えと完全電源断による。

\end{document}
